% =============================================================================
% demo.tex -- a succinct tour of the union-doc.cls feature set.
%
% Compile from this directory (latex-class/demo/v1/) with:
%     lualatex --shell-escape demo.tex
%     lualatex --shell-escape demo.tex    % second pass fills in the ToC
%
% File location matters: the class resolves ../../Font_Assets/ relative to
% this .tex, so demo.tex must live two folders deep inside latex-class/.
% =============================================================================

\documentclass{../../union-doc}

% ----- 1. Cover-page configuration ------------------------------------------
% Every \docXxx hook below is optional; comment any of them out to see the
% cover page degrade gracefully.
\renewcommand{\docPdfTitle}{union-doc.cls Feature Demo}
\renewcommand{\docTitle}{union-doc.cls Feature Demo}
\renewcommand{\docSubtitle}{A one-page tour of every public command}
\renewcommand{\docCoverNote}{Compiled \today\ with LuaLaTeX}
\renewcommand{\docTocDepth}{2}              % 1 = articles only, 2 = + sections

% \DocLogo searches the exposed \docLogoPath (and \docSubtitledLogoPath)
% for you and gracefully falls back to a plain red wordmark if the file is
% missing, so there is no need to write ../../Commonly_Used_Logos/ by hand.
\renewcommand{\docCoverLogo}{\DocLogo{Primary_Logo_Red.png}}

% Optional footer under the TOC (centered italic note).
\renewcommand{\docTocFooter}{%
  \begin{center}\itshape\small
    This document is a non-binding demonstration of \texttt{union-doc.cls}.
  \end{center}%
}

\begin{document}

% ----- 2. Front matter ------------------------------------------------------
% Emits cover logo, title block, cover note, and the boxed table of contents.
\MakeUnionDocFrontMatter

% ----- 3. Articles, sections, and body paragraphs ---------------------------
\DocArticle{art:structure}{Article I --- Structural Commands}
\DocParagraph{\texttt{\textbackslash DocArticle} renders a red, centered
  heading and adds it to the table of contents. Its first argument is an
  internal hyperlink label (used by \texttt{\textbackslash hyperref}), and
  the second is the visible title.}

\DocSection{Section 1 --- Paragraphs and notes}
\DocParagraph{\texttt{\textbackslash DocParagraph} wraps a block of body
  text in the correct font, justification, and clause spacing. Use it
  instead of writing raw text between commands.}

\DocBracketNote{[\texttt{\textbackslash DocBracketNote} renders italic
  parenthetical commentary that is not itself a numbered clause.]}

\DocParagraph{Footnotes are written with
  \texttt{\textbackslash DocFootnote}\DocFootnote{demo}{This is a footnote
  generated by \texttt{\textbackslash DocFootnote\{demo\}\{...\}}.}. The
  first argument is an identifier reserved for future cross-referencing;
  the second is the footnote body.}

% ----- 4. Nested numbered lists ---------------------------------------------
\DocSection{Section 2 --- Nested numbered lists}
\DocParagraph{The class preconfigures four \texttt{enumerate} levels:
  \textbf{1.}, \textbf{i.}, \textbf{a.}, and \textbf{1)}. Just nest the
  standard \texttt{enumerate} environment --- do not load
  \texttt{enumitem} yourself.}

\begin{enumerate}
  \item Top-level clause (\texttt{1.} label).
  \begin{enumerate}
    \item Second-level clause (\texttt{i.} label).
    \begin{enumerate}
      \item Third-level clause (\texttt{a.} label).
      \begin{enumerate}
        \item Fourth-level clause (\texttt{1)} label).
      \end{enumerate}
    \end{enumerate}
  \end{enumerate}
  \item Another top-level clause, to show spacing between items.
\end{enumerate}

% ----- 5. Colors and inline formatting --------------------------------------
\DocSection{Section 3 --- Brand colors}
\DocParagraph{Three named colors are defined:
  \textcolor{UnionRed}{\textbf{UnionRed}} (\texttt{\#D6001C}),
  \textbf{UnionBlack} (\texttt{\#000000}), and
  \colorbox{UnionBlack}{\textcolor{UnionWhite}{\textbf{UnionWhite}}}
  (\texttt{\#FFFFFF}). Use them anywhere with \texttt{\textbackslash color}
  or \texttt{\textbackslash textcolor}.}

% ----- 6. Asset paths and logo helper ---------------------------------------
\DocSection{Section 4 --- Assets, fonts, and logos}
\DocParagraph{The class exposes three overridable asset paths:
  \texttt{\textbackslash docFontPath} (default
  \texttt{../../Font\_Assets/}),
  \texttt{\textbackslash docLogoPath} (default
  \texttt{../../Commonly\_Used\_Logos/}), and
  \texttt{\textbackslash docSubtitledLogoPath} (default
  \texttt{../../Subtitled\_Logos/}). Both logo directories are added to
  \texttt{\textbackslash graphicspath} and \texttt{\textbackslash svgpath}
  at class-load time, so you can reference any logo by its bare filename:}

\begin{center}
  \includegraphics[width=1.2in]{Primary_Logo_Black.png}\qquad
  \DocLogo[width=1.2in]{Primary_Logo_Red.png}
\end{center}

\DocParagraph{The left image above was loaded with a plain
  \texttt{\textbackslash includegraphics\{Primary\_Logo\_Black.png\}}
  --- no path prefix required. The right image uses the
  \texttt{\textbackslash DocLogo[$\langle$opts$\rangle$]\{$\langle$file$\rangle$\}}
  convenience command, which resolves the same search path and falls back
  to a red ``Rensselaer Union'' wordmark if the file is missing.}

\DocParagraph{A monospace family, {\docMonoFont RPI Geist Mono}, is also
  declared on the same \texttt{\textbackslash docFontPath} and is available
  as \texttt{\textbackslash docMonoFont} for inline use.}

\DocSection{Section 5 --- Call-outs and reusable title blocks}
\DocParagraph{\texttt{DocPanel} is a red-bordered tcolorbox environment for
  highlighting a notice or pull quote anywhere in the body:}

\begin{DocPanel}
  This is a \texttt{DocPanel}. Use it to draw attention to an editorial
  note, a ratification stamp, or any block of text that should visually
  stand apart from the surrounding clauses.
\end{DocPanel}

\DocParagraph{\texttt{\textbackslash DocTitleBlock\{title\}\{subtitle\}\{note\}}
  renders a centered hero block. Any argument may be left empty:}

\DocTitleBlock%
  {Section Divider}%
  {Used mid-document to announce a new part}%
  {All three arguments are optional --- leave empty to suppress}

% ----- 6. Amendments --------------------------------------------------------
\DocAmendment{amend:one}{Amendment 1 --- Adding an Amendment Block}
\DocAmendmentSubtitle{Adopted 2026-04-21}{Demonstration subtitle}
\DocParagraph{\texttt{\textbackslash DocAmendment} is visually identical
  to \texttt{\textbackslash DocArticle} but semantically marks an
  amendment. \texttt{\textbackslash DocAmendmentSubtitle} takes a required
  subtitle and an optional secondary note (pass an empty second argument
  to omit it).}

\DocAmendment{amend:two}{Amendment 2 --- An Amendment Without a Note}
\DocAmendmentSubtitle{Adopted 2026-04-21}{}
\DocParagraph{When the second argument of
  \texttt{\textbackslash DocAmendmentSubtitle} is empty, the note line is
  silently suppressed.}

% ----- 7. That is the complete public API -----------------------------------
\DocArticle{art:done}{Article II --- Summary}
\DocParagraph{You have now seen every public command exposed by
  \texttt{union-doc.cls}: the cover hooks
  (\texttt{\textbackslash docPdfTitle}, \texttt{\textbackslash docTitle},
  \texttt{\textbackslash docSubtitle}, \texttt{\textbackslash docCoverNote},
  \texttt{\textbackslash docCoverLogo}, \texttt{\textbackslash docTocDepth},
  \texttt{\textbackslash docTocFooter}), the front-matter builder
  \texttt{\textbackslash MakeUnionDocFrontMatter}, and the body commands
  \texttt{\textbackslash DocArticle}, \texttt{\textbackslash DocSection},
  \texttt{\textbackslash DocParagraph}, \texttt{\textbackslash DocAmendment},
  \texttt{\textbackslash DocAmendmentSubtitle},
  \texttt{\textbackslash DocBracketNote}, and
  \texttt{\textbackslash DocFootnote}.}

\end{document}
